\begin{proofbox}
\textbf{方法一：最快证明（平方和一行法，推荐首先掌握）。}

设 $P(x_0,y_0)\in C$。考虑直线 $\ell:x_0x/a^2+y_0y/b^2=1$；它显然经过 $P$。若 $Q(x,y)$ 同时在 $\ell$ 与椭圆 $C$ 上，则
\[
 \underbrace{\frac{(x-x_0)^2}{a^2}+\frac{(y-y_0)^2}{b^2}}_{\ge0}
 =\underbrace{\left(\frac{x^2}{a^2}+\frac{y^2}{b^2}\right)}_{1}
 +\underbrace{\left(\frac{x_0^2}{a^2}+\frac{y_0^2}{b^2}\right)}_{1}
 -2\underbrace{\left(\frac{x_0x}{a^2}+\frac{y_0y}{b^2}\right)}_{1}=0.
\]
因此 $Q=P$；直线 $\ell$ 与非退化椭圆恰有一个公共点，故 $\ell$ 是 $P$ 处切线。\textbf{无求导、无判别式、无斜率分类，一次覆盖所有切点。}

\medskip
\textbf{方法二：纯几何思路（把椭圆缩放成圆）。}

令 $u=x/a,\ v=y/b$，椭圆变为 $u^2+v^2=1$，切点变为 $U(u_0,v_0)=(x_0/a,y_0/b)$。
单位圆切线与半径垂直，因此
\[
 (u-u_0)u_0+(v-v_0)v_0=0
 \quad\Longrightarrow\quad uu_0+vv_0=u_0^2+v_0^2=1.
\]
代回原坐标即 $x_0x/a^2+y_0y/b^2=1$。可逆线性缩放把圆的切线对应为椭圆的切线（保留相交与相切的关系）。

\medskip
\textbf{方法三：割线逼近切线（不使用导数公式）。}

取椭圆上另一点 $Q(x,y)\ne P$（取 $Q$ 在 $P$ 附近，因此也不会是对径点 $-P$），将 $P,Q$ 的椭圆方程相减：
\[
 \frac{x^2-x_0^2}{a^2}+\frac{y^2-y_0^2}{b^2}=0,
\]
也就是
\[
 \frac{x+x_0}{a^2}(x-x_0)+\frac{y+y_0}{b^2}(y-y_0)=0.
\]
所以割线 $PQ$ 的法向量可取 $((x+x_0)/a^2,(y+y_0)/b^2)$。令 $Q\to P$，法向量趋向 $(2x_0/a^2,2y_0/b^2)$，于是切线满足
\[
 \frac{x_0}{a^2}(X-x_0)+\frac{y_0}{b^2}(Y-y_0)=0.
\]
利用 $P\in C$ 整理得 $x_0X/a^2+y_0Y/b^2=1$。这里 $X,Y$ 表示切线上任意点的坐标，以区别于趋近 $P$ 的动点 $Q$。

\medskip
\textbf{方法四：代数法（隐函数求导，适合求斜率）。}

椭圆两边对 $x$ 求导：
\[
 \frac{2x}{a^2}+\frac{2y}{b^2}y'=0
 \quad\Longrightarrow\quad
 y'=-\frac{b^2x}{a^2y}\qquad(y\ne0).
\]
若 $y_0\ne0$，切线的点斜式为
\[
 y-y_0=-\frac{b^2x_0}{a^2y_0}(x-x_0).
\]
移项并用 $x_0^2/a^2+y_0^2/b^2=1$，即得
\[
 \frac{x_0x}{a^2}+\frac{y_0y}{b^2}=1.
\]
若 $y_0=0$，切点必为 $(\pm a,0)$，竖直切线为 $x=\pm a$，仍满足切点式；原来的斜率公式不能直接代入。

\medskip
\textbf{方法五：不等式法（数形结合：整条椭圆在切线一侧）。}

设 $Q(x,y)$ 是椭圆任意一点。应用柯西不等式：
\[
 \begin{aligned}
 \frac{x_0x}{a^2}+\frac{y_0y}{b^2}
 &\le \sqrt{\frac{x_0^2}{a^2}+\frac{y_0^2}{b^2}}
       \sqrt{\frac{x^2}{a^2}+\frac{y^2}{b^2}}\\
 &=1.
 \end{aligned}
\]
等号要求 $(x/a,y/b)$ 与 $(x_0/a,y_0/b)$ 同向且同长，所以只能有 $Q=P$。因此直线 $x_0x/a^2+y_0y/b^2=1$ 在 $P$ 处接触椭圆，且椭圆整体位于直线的一侧。这直接说明\textbf{“相切”而非“相交穿过”}。

\medskip
\textbf{方法六：参数方程与三角恒等式（把接触条件化成 $\cos0=1$）。}

设切点 $P=(a\cos\theta,b\sin\theta)$，椭圆上的动点 $Q=(a\cos t,b\sin t)$，则
\[
 \frac{x_0x}{a^2}+\frac{y_0y}{b^2}
 =\cos\theta\cos t+\sin\theta\sin t
 =\cos(t-\theta)\le1.
\]
等号当且仅当 $t=\theta+2k\pi$，即 $Q=P$。故直线 $x_0x/a^2+y_0y/b^2=1$ 只在 $P$ 处接触椭圆。它直观地把\textbf{切点}解释为\textbf{角度差为零的唯一位置}。

\medskip
\textbf{推论一：外点的切点弦（极线）证明。}

设 $M(x_1,y_1)$ 在椭圆外，两条实切线的切点为 $T_1(u_1,v_1)$ 和 $T_2(u_2,v_2)$。由已证切点式，$T_i$ 处的切线为
\[
 \frac{u_i x}{a^2}+\frac{v_i y}{b^2}=1\qquad(i=1,2).
\]
两条切线都经过 $M$，故
\[
 \frac{x_1u_i}{a^2}+\frac{y_1v_i}{b^2}=1\qquad(i=1,2).
\]
这说明 $T_1,T_2$ 同时在直线
\[
 \boxed{\frac{x_1x}{a^2}+\frac{y_1y}{b^2}=1}
\]
上，所以它就是切点弦。反过来，这条直线与椭圆的两个交点 $T_i$ 也都满足“$M$ 在 $T_i$ 处切线上”的条件，因此确实是这两条切线的切点。

外点条件不能省略：令 $U=(x_1/a,y_1/b)$，当 $|U|>1$ 时，直线 $u_1u+u_2v=1$（其中 $u_1=x_1/a,u_2=y_1/b$）距单位圆圆心 $1/|U|<1$，确有两个交点；$|U|=1$ 时一个；$0<|U|<1$ 时没有。

\medskip
\textbf{推论二：一般直线的相切判定。}

对椭圆上任意 $Q(x,y)$，柯西不等式给出
\[
 |Ax+By|\le\sqrt{a^2A^2+b^2B^2}.
\]
椭圆上 $Ax+By$ 的最大值和最小值分别为
$\sqrt{a^2A^2+b^2B^2}$ 与 $-\sqrt{a^2A^2+b^2B^2}$，且各在唯一点取得。因此 $Ax+By+C=0$ 与椭圆相切当且仅当
\[
 \boxed{C^2=a^2A^2+b^2B^2},\qquad(A,B)\ne(0,0).
\]
\end{proofbox}
